% !TeX root = ../main.tex
\begin{frame}{02 研究现状——现有研究的不足}
  \begin{SlideCanvas}
    \BodyAt{55}{84}{850}{现有研究已在变形调控与多目标优化方面取得进展。本报告进一步聚焦规定部署路径下的长度控制和径向力响应定制。}
    \PointAt{55}{153}{405}{现有研究仍需解决的问题}
    \PointAt{505}{153}{400}{拟提出的针对性方案}
    \MainRuleAt{190}
    \RedAt{55}{205}{405}{终态指标不约束全过程}
    \DetailAt{55}{239}{405}{终态短缩率接近零，仍不足以判断展开过程中长度是否稳定。}
    \RedAt{505}{205}{400}{全过程长度约束}
    \DetailAt{505}{239}{400}{在规定直径区间限制最大长度偏差，降低病变覆盖长度的不确定性。}
    \RuleAt{289}
    \RedAt{55}{300}{405}{泊松比不充分表征支撑能力}
    \DetailAt{55}{334}{405}{长度响应相近的构型仍可具有不同刚度，支撑需求需单独检验。}
    \RedAt{505}{300}{400}{长度与径向力曲线联合定制}
    \DetailAt{505}{334}{400}{以课题组联合编程为基础，在长度约束下匹配目标径向力响应。}
    \RuleAt{384}
    \RedAt{55}{395}{405}{单胞性能向器件迁移受限}
    \DetailAt{55}{429}{405}{周期单胞不能完整反映曲率、端部与接触作用，需补充径向部署。}
    \RedAt{505}{395}{400}{按径向工况验证可行范围}
    \DetailAt{505}{429}{400}{统一材料、杆宽和应变约束，比较圆管长度、径向力与局部应变。}
    \Sources{研究依据：\RefCarbonaro；\RefWang；\RefYasmin。\\联合编程基础：\RefGroup。右列为拟议支架方案，尚需按径向工况验证。}
  \end{SlideCanvas}
  \BICConclusionBand{拟议方案聚焦\BICBandEmph{长度轨迹与径向承载的共同设计}}
  \note{本页用三组具体问题引出创新思路。已有负／零泊松比器件、多目标优化和全过程数值模型，不能将零泊松比或共同考虑短缩与径向力称为首次提出。第一行强调指标的逻辑范围：单个终态短缩率不能充分确定整个路径；不据此断言所有现有论文都只考察终态。Wang2021已有原始、最大扩张和回弹状态的长度统计，Yasmin2026也包含压握与释放及多目标设计。第二行依据Carbonaro2024的等效刚度—器件性能映射，说明泊松比或长度耦合不足以独立确定支撑性能；课题组方法提供几何与非线性载荷响应共同编程的基础，拟进一步扩展至指定径向加载分支内的目标曲线和可行范围。第三行指出周期单胞与有限圆管、接触部署之间的模型差别，拟在统一材料和制造约束下进行器件对照。原课题组圆管结果属于轴向加载，尚不是径向支撑证据。右列均为待验证的支架应用构想，长度控制有助于覆盖长度稳定，仍需检查端部位置、接触和相应医学指标。文献doi：10.1016/j.cmpb.2024.108467；10.1016/j.eng.2020.02.013；10.3390/mi17060736；10.1016/j.tws.2026.114956。}
\end{frame}
